\section*{الفصل \ref{ch:g10:electron-shells} --- الطبقات الإلكترونية والجدول الدوري}

\begin{solution}{exo:g10:electron-shells:1}
C: $1\mathrm{s}^2\,2\mathrm{s}^2\,2\mathrm{p}^2$. N:
$1\mathrm{s}^2\,2\mathrm{s}^2\,2\mathrm{p}^3$. Mg:
$1\mathrm{s}^2\,2\mathrm{s}^2\,2\mathrm{p}^6\,3\mathrm{s}^2$.
\end{solution}

\begin{solution}{exo:g10:electron-shells:2}
الكربون 4، والنيتروجين 5، والمغنيسيوم 2.
\end{solution}

\begin{solution}{exo:g10:electron-shells:3}
\omterm{def:g10:electron-shells:configuration}{الطبقة} الخارجية $n = 3$: \omterm{def:g9:periodic-table-first-look:period}{الدورة} 3. و$2 + 3 = 5$ \omterm{def:g9:inside-the-atom:nucleus}{إلكترونات} تكافؤ: \omterm{def:g9:periodic-table-first-look:period}{الزمرة} 15
(الفوسفور).
\end{solution}

\begin{solution}{exo:g10:electron-shells:4}
\omterm{def:g10:electron-shells:family}{الفلزات القلوية} (الليثيوم، والصوديوم، والبوتاسيوم)؛ \omterm{def:g10:electron-shells:family}{والهالوجينات} (الفلور،
والكلور، والبروم، واليود)؛ \omterm{def:g10:electron-shells:family}{والغازات النبيلة} (الهيليوم، والنيون، والأرغون).
\end{solution}

\begin{solution}{exo:g10:electron-shells:5}
تميل \omterm{def:g7:atoms-and-molecules:atom}{الذرات} إلى اكتساب \omterm{def:g9:inside-the-atom:nucleus}{إلكترونات} أو فقدها لتصير في طبقتها الخارجية ثمانية
\omterm{def:g9:inside-the-atom:nucleus}{إلكترونات}، مثل أقرب \omterm{def:g10:electron-shells:family}{غاز نبيل}.
\end{solution}

\begin{solution}{exo:g10:electron-shells:6}
K، $2.8.8.1$، يفقد واحدًا: \ce{K+}،
$1\mathrm{s}^2\,2\mathrm{s}^2\,2\mathrm{p}^6\,3\mathrm{s}^2\,3\mathrm{p}^6$
(الأرغون). F، $2.7$، يكتسب واحدًا: \ce{F-}،
$1\mathrm{s}^2\,2\mathrm{s}^2\,2\mathrm{p}^6$ (النيون).
\end{solution}

\begin{solution}{exo:g10:electron-shells:7}
$1\mathrm{s}^2\,2\mathrm{s}^2\,2\mathrm{p}^6\,3\mathrm{s}^2$: $Z = 12$،
المغنيسيوم.
\end{solution}

\begin{solution}{exo:g10:electron-shells:8}
\ce{Mg^{2+}} و\ce{Cl-}: \ce{MgCl2}. \ce{Al^{3+}} و\ce{O^{2-}}:
\ce{Al2O3}.
\end{solution}

\begin{solution}{exo:g10:electron-shells:9}
الكبريت (6 \omterm{def:g9:inside-the-atom:nucleus}{إلكترونات} تكافؤ، $2.8.6$)، في \omterm{def:g9:periodic-table-first-look:period}{الزمرة} نفسها، 16.
\end{solution}

\begin{solution}{exo:g10:electron-shells:10}
طبقته الخارجية ممتلئة أصلًا (ثمانية \omterm{def:g9:inside-the-atom:nucleus}{إلكترونات}): فلا شيء يكسبه بفقد \omterm{def:g9:inside-the-atom:nucleus}{إلكترونات}
أو اكتسابها.
\end{solution}

\begin{solution}{exo:g10:electron-shells:11}
للأرغون 18 إلكترونًا؛ و\ce{X^{2+}} فقد 2، فلذرة X عشرون إلكترونًا: إنه الكالسيوم.
\end{solution}

\begin{solution}{exo:g10:electron-shells:12}
الهيليوم، $1\mathrm{s}^2$، طبقته الخارجية ممتلئة (\omterm{def:g10:electron-shells:configuration}{فالطبقة} 1 لا تتسع إلا
لإلكترونين)، مثل \omterm{def:g10:electron-shells:family}{الغازات النبيلة} الأخرى؛ ولا يسلك سلوك \omterm{def:g9:periodic-table-first-look:metal}{فلزات} \omterm{def:g9:periodic-table-first-look:period}{الزمرة}
2.
\end{solution}

\begin{solution}{exo:g10:electron-shells:13}
\omterm{def:g9:inside-the-atom:nucleus}{إلكترون} تكافئه في \omterm{def:g10:electron-shells:configuration}{الطبقة} 4، أبعد عن \omterm{def:g9:inside-the-atom:nucleus}{النواة} من \omterm{def:g9:inside-the-atom:nucleus}{إلكترون} الصوديوم في \omterm{def:g10:electron-shells:configuration}{الطبقة} 3:
فهو أضعف ارتباطًا، ويُفقد بسهولة أكبر.
\end{solution}

\begin{solution}{exo:g10:electron-shells:14}
لكل \omterm{def:g9:ions:ion}{الأيونات} الأربعة 18 إلكترونًا \omterm{def:g10:electron-shells:configuration}{والتوزيع الإلكتروني} للأرغون:
\[ 1\mathrm{s}^2\,2\mathrm{s}^2\,2\mathrm{p}^6\,3\mathrm{s}^2\,3\mathrm{p}^6. \]
\end{solution}

\begin{solution}{exo:g10:electron-shells:15}
يكتسب 3 \omterm{def:g9:inside-the-atom:nucleus}{إلكترونات} ليبلغ الثمانية: فله 5 \omterm{def:g9:inside-the-atom:nucleus}{إلكترونات} تكافؤ، \omterm{def:g9:periodic-table-first-look:period}{الزمرة}
15. وفي \omterm{def:g9:periodic-table-first-look:period}{الدورة} 3 هو الفوسفور:
\[ 1\mathrm{s}^2\,2\mathrm{s}^2\,2\mathrm{p}^6\,3\mathrm{s}^2\,3\mathrm{p}^3. \]
\end{solution}

\begin{solution}{pb:g10:electron-shells:1}
\textbf{1.} الألومنيوم: 13 بروتونًا، و13 إلكترونًا. الأكسجين: 8 \omterm{def:g9:inside-the-atom:proton}{بروتونات}، و8
\omterm{def:g9:inside-the-atom:nucleus}{إلكترونات}.

\textbf{2.} Al: $1\mathrm{s}^2\,2\mathrm{s}^2\,2\mathrm{p}^6\,3\mathrm{s}^2\,3\mathrm{p}^1$.
O: $1\mathrm{s}^2\,2\mathrm{s}^2\,2\mathrm{p}^4$.

\textbf{3.} Al: 3 \omterm{def:g9:inside-the-atom:nucleus}{إلكترونات} تكافؤ، \omterm{def:g9:periodic-table-first-look:period}{الدورة} 3، \omterm{def:g9:periodic-table-first-look:period}{الزمرة} 13. O: 6، \omterm{def:g9:periodic-table-first-look:period}{الدورة} 2،
\omterm{def:g9:periodic-table-first-look:period}{الزمرة} 16.

\textbf{4.} الألومنيوم \omterm{def:g9:periodic-table-first-look:metal}{فلز}، والأكسجين \omterm{def:g9:periodic-table-first-look:metal}{لافلز}.

\textbf{5.} \ce{Al^{3+}}: $1\mathrm{s}^2\,2\mathrm{s}^2\,2\mathrm{p}^6$.

\textbf{6.} \ce{O^{2-}}: $1\mathrm{s}^2\,2\mathrm{s}^2\,2\mathrm{p}^6$.

\textbf{7.} النيون.

\textbf{8.} الألومنيوم يفقد 3، والأكسجين يكتسب 2.

\textbf{9.} $2 \times (+3) + 3 \times (-2) = 0$: \ce{Al2O3}.

\textbf{10.} المفقودة: $2 \times 3 = 6$. المكتسبة: $3 \times 2 = 6$.

\textbf{11.} \omterm{def:g9:inside-the-atom:nucleus}{الإلكترونات} لا تُخلق ولا تفنى: فالتي فقدها الألومنيوم هي نفسها
التي اكتسبها الأكسجين، والمركب متعادل.

\textbf{12.} يحمل الشحنة نفسها، $+3$: فتبقى شحنات البلورة
متوازنة.

\textbf{13.} $2 \times 27 + 3 \times 16 = 102$ نوكليونًا.

\textbf{14.} $102 \times 1.67 \times 10^{-27} = \qty{1.70e-25}{kg}$.

\textbf{15.} $10^{-3} / 1.70 \times 10^{-25} \approx 5.87 \times 10^{21}$
وحدة صيغة.

\textbf{16.} $2 \times 5.87 \times 10^{21} = 1.17 \times 10^{22}$ \omterm{def:g9:ions:ion}{أيون}
\ce{Al^{3+}}؛ و$3 \times 5.87 \times 10^{21} = 1.76 \times 10^{22}$ \omterm{def:g9:ions:ion}{أيون}
\ce{O^{2-}}.

\textbf{17.} \omterm{def:g9:inside-the-atom:nucleus}{الإلكترون} أخف من \omterm{def:g9:inside-the-atom:proton}{النوكليون} بنحو 1836 مرة؛ \omterm{def:g9:inside-the-atom:nucleus}{والإلكترونات} الخمسون
في وحدة الصيغة تزن أقل من ثلاثة أجزاء من مئة من \omterm{def:g9:inside-the-atom:proton}{نوكليون}
واحد.

\textbf{18.} $6 \times 5.87 \times 10^{21} \approx 3.5 \times 10^{22}$
\omterm{def:g9:inside-the-atom:nucleus}{إلكترون} منقول لصنع \qty{1}{g} من الكوراندوم.
\end{solution}
